\chapter*{Аннотация}
\addcontentsline{toc}{chapter}{Аннотация}

В диссертационной работе рассмотрены модели нелокальной теплопроводности и термоупругости. В рассматриваемых моделях, уравнения теплопроводности и равновесия не изменили свою формулировку по сравнению с классическими моделями механики сплошных сред, но представлены в интегро-дифференциальном виде за счёт применения к гипотезам Био – Фурье и Дюамеля – Неймана интегрального нелокального оператора. Для поиска решений был разработан численный метод на основе метода конечных элементов с адаптацией под многопроцессорные вычислительные системы. Помимо этого, в работе рассмотрены различные варианты предобуславливания систем линейных алгебраических уравнений, которые получены после дискретизации уравнений. В рамках работы над диссертацией реализован авторский программный комплекс NonLocFEM, при написании которого была поставлена задача эффективной реализации предложенных алгоритмов. Исследована применимость принципа Сен-Венана к задачам нелокальной упругости, а также его аналога для тепловых задач – принципа стабильности тепловых потоков. Полученные решения свидетельствуют о наличии кромочного эффекта, который характеризуется снижением уровня напряжения и плотности теплового потока на свободных от условий границах и повышении температуры и модуля вектора перемещения на границах с нагружением. Помимо этого, были исследованы области с концентраторами, на примере задач о растяжении Т-образной пластины, задачи Кирша с обобщением на эллиптические вырезы и задачи о проходящем тепловом потоке сквозь пластину с вырезом. Продемонстрированные в работе решения свидетельствуют о снижении роли концентраторов и повышении уровня напряжения и плотности тепловых потоков внутри рассматриваемых областей.